\SourceLeaf{394}%
woraus nach den gehörigen Entwickelungen leicht folgt
\[
\begin{aligned}
&4\theta\sin(A''-A')\sin(A-A'')\sin(A'-A)\\
&\qquad=b\sin2(A''-A')+b'\sin2(A-A'')+b''\sin2(A'-A)
\end{aligned}
\]
\[
\begin{aligned}
&4uu\sin(A''-A')^2\sin(A-A'')^2\sin(A'-A)^2\\
&\qquad=bb\sin(A''-A')^2\\
&\qquad\quad+b'b'\sin(A-A'')^2\\
&\qquad\quad+b''b''\sin(A'-A)^2\\
&\qquad\quad+2b'b''\cos(A''-A')\sin(A-A'')\sin(A'-A)\\
&\qquad\quad+2bb''\sin(A''-A')\cos(A-A'')\sin(A'-A)\\
&\qquad\quad+2bb'\sin(A''-A')\sin(A-A'')\cos(A'-A)
\end{aligned}
\]
und hieraus
\[
\begin{aligned}
&4(\theta\theta-uu)\sin(A''-A')^2\sin(A-A'')^2\sin(A'-A)^2\\
&\qquad=-bb\sin(A''-A')^4\\
&\qquad\quad-b'b'\sin(A-A'')^4\\
&\qquad\quad-b''b''\sin(A'-A)^4\\
&\qquad\quad+2b'b''\sin(A-A'')^2\sin(A'-A)^2\\
&\qquad\quad+2bb''\sin(A''-A')^2\sin(A'-A)^2\\
&\qquad\quad+2bb'\sin(A''-A')^2\sin(A''-A)^2
\end{aligned}
\]
\par
Ich habe diese Formeln hierher gesetzt, weil sie auch in andern Fällen zuweilen mit Nutzen zu gebrauchen sind. Man sieht leicht, dass das, was auf der rechten Seite steht, das Product aus dem vier Factoren sei.
\[
\begin{array}{l}
+\BareRoot{b}.\sin(A''-A')+\BareRoot{b'}.\sin(A-A'')+\BareRoot{b''}.\sin(A'-A)\\
-\BareRoot{b}.\sin(A''-A')+\BareRoot{b'}.\sin(A-A'')+\BareRoot{b''}.\sin(A'-A)\\
+\BareRoot{b}.\sin(A''-A')-\BareRoot{b'}.\sin(A-A'')+\BareRoot{b''}.\sin(A'-A)\\
+\BareRoot{b}.\sin(A''-A')+\BareRoot{b'}.\sin(A-A'')-\BareRoot{b''}.\sin(A'-A)
\end{array}
\]
\par
Substituirt man hier für \(b,b',b''\) ihre Werthe und setzt Kürze halber
\[
a\sin(A''-A')+a'\sin(A-A'')+a''\sin(A'-A)=M
\]
so wird
\[
(\theta\theta-uu)\sin(A''-A')^2\sin(A-A'')^2\sin(A'-A)^2
\]
gleich dem Producte aus den vier Factoren
